El comportamiento observado evidencia que el crecimiento de las remesas familiares ha estado acompañado por una expansión del Producto Interno Bruto (PIB) durante el período 2015–2025. A medida que aumenta el ingreso de divisas provenientes del exterior, también se incrementa la capacidad de consumo e inversión de los hogares receptores, lo que contribuye al dinamismo de distintos sectores de la economía y favorece el crecimiento de la producción nacional.
Aunque a partir de 2019 la dispersión de los datos es ligeramente mayor, principalmente por los efectos económicos derivados de la pandemia de COVID-19, la trayectoria general continúa mostrando un comportamiento creciente. Esto indica que, aun bajo escenarios de incertidumbre, las remesas mantuvieron un papel importante como fuente de ingresos para los hogares guatemaltecos y como mecanismo de apoyo para la recuperación de la actividad económica.
Desde una perspectiva macroeconómica, estos resultados resaltan la relevancia de las remesas dentro del modelo económico guatemalteco. Sin embargo, también evidencian una dependencia significativa de recursos provenientes del exterior, los cuales pueden verse afectados por cambios en las condiciones económicas internacionales, el mercado laboral de los países donde residen los migrantes o modificaciones en las políticas migratorias. Por ello, resulta conveniente que las autoridades económicas complementen este flujo de ingresos con políticas orientadas a fortalecer la inversión, la productividad y la generación de empleo interno, reduciendo gradualmente la dependencia de factores externos como fuente de crecimiento económico.